% =============================================================
% Chapitre 5 --- Réalisation et industrialisation (~8 pages)
% Source : ../../CH5_Realisation.md
% =============================================================
\chapter{Réalisation et industrialisation}
\label{ch:realisation}
\soustitrechapitre{Description en code, chaînes de livraison et de provisionnement, supervision
et accueil d'un locataire}

% Nomenclature des schémas : ligne 1 = rôle fonctionnel, ligne 2 = produit
% nommé, en gris et au plancher typographique de 8 pt (\scriptsize).
% \gcprod est désormais défini dans config/convention_graphique.tex, plusieurs
% chapitres l'employant.

Le chapitre précédent décrivait une architecture cible ; celui-ci rend compte de ce qui a été
construit, en distinguant partout \textbf{réalisé et vérifié}, \textbf{réalisé et rattaché à une
vérification planifiée} et \textbf{écarté par une décision motivée}. Deux chaînes structurent
l'exposé : la livraison applicative et le provisionnement de l'infrastructure.

\section{Environnement de travail et description en code}
\label{sec:env-travail}

Deux environnements réels correspondent à deux projets cloud distincts : le développement,
dimensionné au minimum, et la recette, de topologie identique à la cible ; la production est une
cible documentée, non provisionnée. Séparer par projet plutôt que par étiquette rend l'isolation
\textbf{structurelle} --- les ressources d'un environnement ne sont pas adressables depuis
l'autre --- sous réserve des liaisons d'autorisation accordées à l'identité de déploiement,
qu'aucun test dédié n'a éprouvées. Le développement n'est pas un miroir fidèle de la recette : la
configuration du second locataire et huit variables d'accueil de plusieurs applications n'existent
qu'en recette, et la chaîne de promotion n'a pas été déroulée de bout en bout. Une règle vaut
cependant sans exception --- \emph{les mêmes modules servent partout, seules les variables de
dimensionnement changent} --- de sorte qu'aucune configuration « sans authentification en
développement » ne risque d'être promue par inadvertance. Le quatrième projet d'amorçage prévu à la
conception n'a pas été créé, ses éléments étant replacés dans chaque projet d'environnement :
simplification cohérente avec le principe de justification (PR1), enregistrée comme écart.

Le dépôt d'infrastructure transpose le modèle en couches du \cref{ch:conception} en un module par couche
(annexe~\ref{ann:extraits-realisation}). Les trois politiques d'organisation prévues --- pas
d'adresse publique de base de données, pas de clé d'identité de service, pas d'espace de stockage
public --- se sont révélées inapplicables dans le périmètre administratif disponible (écart É7) :
ces interdictions sont donc recherchées par la configuration des modules, les contrôles de la
chaîne d'infrastructure et la revue du \emph{plan}.

La \cref{fig:capture-K36} restitue l'inventaire de ces charges tel que la plateforme le publie,
à mettre en regard de l'inventaire technologique du \cref{ch:etat-art}.

\begin{figure}[!htbp]
\centering
\IfFileExists{figures/K36.png}{%
  \setlength{\fboxrule}{0.4pt}\setlength{\fboxsep}{0pt}%
  \fcolorbox{black!30}{white}{\includegraphics[width=0.86\textwidth]{figures/K36.png}}%
}{%
  \begin{minipage}{0.86\textwidth}
  \begin{extraitconf}[breakable=false, fontupper=\footnotesize,
      title={Cadre de capture K36 : spécification de prise}]
  \textbf{Vue :} liste des services de la plateforme d'exécution de conteneurs, projet de recette, colonnes « service », « région » et « dernier déploiement » visibles.\par\smallskip
  \textbf{Date de prise :} sur le projet de recette. \emph{La date portée par la capture fait foi.}\par\smallskip
  \textbf{Contenu probant :} les charges de travail du socle nommées et actives dans une région européenne unique, ce qui établit d'un seul coup deux propriétés énoncées au \cref{ch:conception} : la séparation par service --- donc une identité par charge --- et l'absence de toute charge hors de la région déclarée. Le décompte visible doit coïncider avec celui de l'inventaire technologique.\par\smallskip
  \textbf{Masquage} (rectangle opaque, jamais de flou) \textbf{:} numéro et identifiant du projet cloud, adresses de service publiées.
  \end{extraitconf}
  \end{minipage}%
}
\caption[Charges de travail du socle sur la plateforme d'exécution]{Les charges de travail du socle sur la plateforme d'exécution de conteneurs (Cloud Run), projet de recette. Restitution de l'inventaire technologique (\cref{sec:trente-trois-briques}) ; la date de relevé est portée par la capture.}
\label{fig:capture-K36}
\end{figure}

\subsection{Frontière entre le code et l'exploitation}
\label{subsec:frontiere-code}

Le principe de description en code (PR4) s'applique à toutes les ressources durables du socle,
\textbf{à six exceptions que ce mémoire nomme et justifie} : trois tiennent à la nature de l'objet,
trois à une limite du dispositif et sont traitées comme des écarts, avec leur contrôle
compensatoire. Les énoncer vaut mieux que revendiquer une couverture totale, car \textbf{un principe
dont on ne sait pas dire la frontière n'est pas appliqué, il est proclamé}.

\begin{enumerate}\itemsep0pt
\item \textbf{Les valeurs des secrets et les données applicatives.} Le code crée les conteneurs,
les clés de chiffrement et les droits d'accès ; jamais le contenu, qu'une description en code ferait
mécaniquement apparaître en clair dans le fichier d'état.
\item \textbf{Les comptes du tableau de bord et leurs secrets de second facteur}, nés d'un appel
authentifié à l'API de supervision à l'exécution ; le premier compte d'administration est créé par
une tâche éphémère plutôt que par une ressource déclarée.
\item \textbf{Les enregistrements de noms de domaine}, tenus chez le bureau d'enregistrement, hors
du périmètre d'administration du fournisseur cloud : dépendance externe au \cref{ch:conception}, non composant
de la couche de périmètre.
\item \textbf{Le versionnement de l'espace de stockage de l'état d'infrastructure}, activé à la
main. Contrôle compensatoire : la revue d'inventaire de cet espace ; condition de clôture : sa
déclaration dans un module d'amorçage distinct de l'état qu'il héberge.
\item \textbf{Les ordres de révocation SQL qui séparent les locataires} sur l'instance partagée
(annexe~\ref{ann:extraits-realisation}) : gestes d'exploitation réversibles par une restauration à
un instant donné, qu'un contrôle quotidien \emph{détecte} sans les \emph{prévenir}
(\cref{sec:locataire}).
\item \textbf{La liaison d'usurpation d'identité accordée à l'opérateur humain} sur l'identité de la
tâche d'enrichissement, posée à la main pour rejouer les tests de bout en bout. Contrôle
compensatoire : elle est nominative et figure à l'inventaire des chaînes de délégation du
\cref{ch:conception}.
\end{enumerate}

S'y ajoute une zone qui n'est pas une exception mais un \emph{partage}, couture la plus délicate
entre les deux chaînes et commune à toute équipe combinant infrastructure en code
\cite{iac_morris2020} et livraison continue : la chaîne d'infrastructure (F7) crée le service
d'exécution et en possède la définition, la chaîne applicative (F3) met à jour l'image à chaque
livraison, et si les deux prétendent posséder le champ « image », chaque exécution de l'une annule
le travail de l'autre. \textbf{Résolution retenue :} l'infrastructure possède le service,
l'applicatif possède l'image, le code d'infrastructure ignorant explicitement ce champ. Le fichier
de variables de la recette n'est par ailleurs pas versionné alors qu'il porte des décisions de
sécurité --- chemins protégés par le filtrage applicatif, services ingérés dans l'entrepôt --- ce
qui limite la reproductibilité d'un provisionnement depuis un poste neuf : c'est le correctif de
reproductibilité le plus rentable identifié.

\subsection{État d'infrastructure : sécurité et détection de dérive}
\label{subsec:derive}

Le fichier d'état décrit la topologie complète du système : donnée sensible, traitée comme telle par
cinq mesures détaillées en annexe~\ref{ann:extraits-realisation}. C'est la règle des secrets
ci-dessus qui l'empêche de devenir un coffre : aucune valeur sensible n'y a été relevée lors des
revues du 19/08/2026, et son inspection exhaustive sous une identité non autorisée est l'objet du
protocole T17, \textbf{conforme} le 20/08/2026 sur quatre refus opposés à trois identités de service
(\cref{sec:campagne}). La détection de dérive \cite{hashicorp_drift} repose sur une planification
régulière : pour une \emph{ressource déjà gérée dans l'état}, un écart entre la configuration
déclarée et la valeur observée produit un plan non vide ; une ressource nouvelle, absente de l'état,
y échappe et exige un contrôle d'inventaire complémentaire. Le protocole T16 modifie hors code un
attribut réversible d'une ressource gérée, planifie avec rafraîchissement, puis rapproche le
résultat des journaux d'audit : il est porté comme partiellement conforme au \cref{ch:validation}.

\begin{alertbox}[boxsep=1pt, left=4pt, right=4pt, top=3pt, bottom=3pt]
Deux limites sont rapportées bien que défavorables. \textbf{Le verrouillage des versions a dû être
corrigé} : le fichier de verrouillage et l'initialisation complète de la validation ont été ajoutés
le 19/08/2026 après constat d'un défaut de reproductibilité. \textbf{Le plan ne prédit pas toujours
l'effet réel de l'application} : la planification ne capture pas certaines validations effectuées
par le fournisseur au moment de l'application, comme le montre l'incident détaillé en
\cref{subsec:incidents} ; la procédure a été amendée en conséquence.
\end{alertbox}

Une fois ces correctifs appliqués, la planification du 19/08/2026 sur la recette ne renvoie
\textbf{aucune différence} : la configuration déclarée dans le dépôt correspond exactement à l'état
réellement provisionné.

\section{Chaîne de livraison applicative}
\label{sec:chaine-livraison}

\subsection{Déclenchement, permissions, graphe des tâches et portes bloquantes}

La chaîne applicative se déclenche sur la poussée sur la branche principale et sur l'ouverture
d'une demande de fusion vers cette branche, \textbf{filtrée par chemins}
(annexe~\ref{ann:extraits-realisation}). Le filtre n'est pas un confort : avant sa pose,
\textbf{toute poussée sur la branche principale reconstruisait et redéployait la plateforme}, y
compris une correction de documentation --- un déclencheur non filtré dilue la signification d'un
déploiement et rend illisible le lien entre un changement et son effet. Les permissions du jeton de
la forge sont, elles, \textbf{déclarées au niveau du flux} --- lecture du dépôt, émission d'un
jeton d'identité --- donc héritées par les six tâches, y compris celles d'analyse et de test qui ne
s'authentifient jamais auprès du fournisseur cloud : la lisibilité d'un fichier unique l'emporte
ici sur la granularité du moindre privilège. Le périmètre reste borné, aucun droit d'écriture sur
le dépôt, mais l'écart est réel ; sa condition de clôture est la factorisation des deux chaînes en
un flux appelable commun, sur le modèle de la chaîne dédiée à l'application pilote, qui applique
déjà le réglage fin.

Les six tâches ne s'exécutent pas en séquence : elles forment un graphe orienté à \textbf{deux
chemins parallèles} qui se rejoignent, celui de l'API et celui du tableau de bord, aux dépendances
\emph{bloquantes}. Vue par étapes, la chaîne de l'API en compte huit, dont trois \textbf{portes
bloquantes} \cite{owasp_cicd_top10} ; les tests de l'API en forment une quatrième, parallèle à
l'analyse, et le chemin du tableau de bord porte la sienne
(\cref{fig:sequence-livraison}). Deux réglages absents complètent ce portrait, assumés plutôt que
tus : la chaîne ne déclare ni contrôle de concurrence --- deux poussées rapprochées se déploient en
parallèle, sans garantie d'ordre --- ni environnement protégé, donc aucune approbation manuelle
avant le déploiement. Les deux relèvent du même arbitrage, une chaîne à un seul opérateur sur un
environnement de recette, et de la même condition de bascule : l'ouverture de la branche principale
à un second contributeur.

\begin{figure}[htbp]
\centering
\begin{tikzpicture}[
  x=1cm, y=1cm, font=\footnotesize,
  etape/.style={nsysteme, text width=1.64cm, inner sep=1.5pt,
                minimum height=2.05cm},
  portec/.style={etape, line width=1.2pt, fill=menalplanB},
  bloc/.style={nsysteme, text width=2.55cm, inner sep=3pt,
               minimum height=0.85cm},
  portel/.style={bloc, line width=1.2pt, fill=menalplanB},
  arret/.style={nabsent, text width=1.45cm, inner sep=2.5pt,
                minimum height=0.50cm, font=\scriptsize},
  verif/.style={nabsent, text width=2.55cm, inner sep=3pt,
                minimum height=0.85cm}]

% ============ chemin de l'API : huit etapes, dont trois portes ==============
\node[etape]  (s1) at (0.90,3.60) {\textbf{1.}\\identité fédérée\\\gcprod{WIF}};
\node[portec] (s2) at (2.78,3.60) {\textbf{2.}\\recherche de secrets\\\gcprod{Gitleaks}};
\node[portec] (s3) at (4.66,3.60) {\textbf{3.}\\analyse statique\\\gcprod{Semgrep}};
\node[etape]  (s4) at (6.54,3.60) {\textbf{4.}\\cons\-truc\-tion de l'image\\\gcprod{Docker}};
\node[portec] (s5) at (8.42,3.60) {\textbf{5.}\\analyse de vulné\-ra\-bi\-lité\\\gcprod{Trivy}};
\node[etape]  (s6) at (10.30,3.60) {\textbf{6.}\\pu\-bli\-ca\-tion au registre\\\gcprod{Artifact Registry}};
\node[etape]  (s7) at (12.18,3.60) {\textbf{7.}\\dé\-ploie\-ment par étiquette\\\gcprod{Cloud Run}};
\node[etape]  (s8) at (14.06,3.60) {\textbf{8.}\\boucle F6\\\gcprod{BigQuery}};
\foreach \a/\b in {s1/s2,s2/s3,s3/s4,s4/s5,s5/s6,s6/s7,s7/s8}%
  {\draw[fluxdonnee] (\a) -- (\b);}

% ============ la porte de tests, parallele, qui bloque la construction ======
\node[portel] (tst) at (6.54,5.65) {porte : tests\\\gcprod{31 tests Python}};
\draw[fluxdonnee] (tst.south) -- (s4.north);

% ============ les refus : la chaine s'arrete ================================
\node[arret] (ar1) at (2.55,5.65) {chaîne\\arrêtée};
\node[arret] (ar2) at (10.05,5.65) {chaîne\\arrêtée};
\draw[fluxrefuse] (s2.north) -- node[croixrefus]{} (ar1.south west);
\draw[fluxrefuse] (s3.north) -- node[croixrefus]{} (ar1.south east);
\draw[fluxrefuse] (s5.north) -- node[croixrefus]{} (ar2.south west);
\draw[fluxrefuse] (tst.east) -- node[croixrefus]{} (ar2.west);

% ============ chemin du tableau de bord =====================================
\node[portel] (t1) at (4.35,1.30) {porte : tests du tableau de bord};
\node[bloc]   (t2) at (8.30,1.30) {construction et déploiement du tableau de bord};
\draw[fluxdonnee] (t1) -- (t2);
\node[arret] (ar3) at (1.30,1.30) {chaîne\\arrêtée};
\draw[fluxrefuse] (t1.west) -- node[croixrefus]{} (ar3.east);
\draw[fluxdonnee] (s1.south) -- (0.90,2.22) -- (4.35,2.22) -- (t1.north);

% ============ la verification, posterieure au deploiement ===================
\node[verif] (ver) at (12.60,1.30) {vérification de bout en bout};
\draw[fluxevenement] (s8.south) |- (ver.east);
\draw[fluxevenement] (t2.east) -- (ver.west);

% ============ bandes de chemin ==============================================
\begin{scope}[on background layer]
  \node[bandedonnee, fit=(s1)(s2)(s3)(s4)(s5)(s6)(s7)(s8)] (bA) {};
  \node[bandedonnee, fit=(t1)(t2)] (bB) {};
\end{scope}
\node[titrebande, anchor=south east] at (bA.north east) {chemin de l'API};
\node[titrebande, anchor=north east] at (bB.south east) {chemin du tableau de bord};

% ============ legende interne ===============================================
\node[legendecadre, text width=14.55cm, anchor=north west, font=\scriptsize]
     at (0.05,0.08)
  {\legendetrait{fluxdonnee}{enchaînement bloquant --- l'étape suivante n'est atteinte que si
   la précédente réussit}\hfill
   \legendetrait{fluxevenement}{non bloquant}\\[2pt]
   \legendecroix{refus : la chaîne s'arrête}\hfill
   cadre épais et fond grisé : \textbf{porte bloquante}\hfill
   cadre en tirets : hors du chemin bloquant};

\end{tikzpicture}
\caption[Chaîne de livraison applicative : étapes, portes bloquantes et refus]{Chaîne de
livraison applicative. Chemin de l'API : huit étapes, dont trois portes bloquantes, plus la porte
de tests qui conditionne la construction ; chemin du tableau de bord : sa propre porte. La
vérification de bout en bout suit le déploiement et ne le défait pas --- c'est une vérification,
non une porte. Les cas de refus sont décrits en \cref{subsec:refus} ; deux d'entre eux portent un
protocole numéroté, T8 et T10, dont le \cref{ch:validation} donne le verdict daté.}
\label{fig:sequence-livraison}
\end{figure}

L'ordre des portes n'est pas arbitraire : la recherche de secrets \cite{gitguardian_sprawl} précède
l'analyse statique pour une raison économique, un secret étant compromis dès sa publication et sa
détection à la première étape déclenchant immédiatement la rotation. Le déploiement à l'étape 7
s'effectue \textbf{par étiquette SHA} plutôt que par tag mutable, dans l'esprit des cadres de
provenance \cite{slsa_framework}, et l'étape 5 s'appuie sur une analyse d'image alignée sur les
pratiques recommandées pour les conteneurs \cite{gcp_container_best_practices}.

\subsection{Les huit familles de contrôles}
\label{subsec:familles-controles}

Décrire une chaîne par ses étapes ne dit pas ce qu'elle contrôle. Lue par \emph{familles de
contrôles}, elle en compte huit, dont \textbf{cinq sont réalisées et bloquantes} --- recherche de
secrets versionnés, analyse statique du code, analyse de la description d'infrastructure, analyse
de vulnérabilité des images, tests automatisés --- et \textbf{trois sont arbitrées} : analyse de
composition des dépendances, test dynamique du service déployé, signature d'artefact et
nomenclature logicielle (décision D07). Chaque arbitrage porte le chemin par lequel le risque est
capté en aval et sa condition de réévaluation. Le \cref{tab:familles-controles} donne les huit,
famille par famille ; sa dernière colonne est celle qu'il faut lire, car la destination du résultat
détermine si un constat est exploitable ou seulement affiché --- et \textbf{une seule famille voit
son résultat quitter la chaîne} : l'analyse de vulnérabilité des images, dont les constats sont
versés à l'entrepôt de supervision.


\textbf{La boucle qui ferme le circuit.} La plupart des chaînes de ce niveau s'arrêtent au constat,
et le rapport meurt dans un journal que personne ne relit ; ici, le résultat de l'analyse d'image
devient une donnée de supervision, interrogeable avec les mêmes outils que les détections, ce qui
rend possible le flux F6 (\cref{subsec:boucle-f6}). La contrepartie doit être dite : aucun constat
des autres familles n'est publié au format d'échange normalisé qu'attend l'onglet de sécurité de la
forge, et les vulnérabilités du tableau de bord ne sont pas exportées --- la chaîne ferme
\emph{une} boucle, pas toutes.

\textbf{La discipline des exclusions d'analyse statique.} Chacune des sept exclusions porte une
justification écrite ancrée sur un fichier et une ligne (annexe~\ref{ann:extraits-realisation}).
Le point remarquable est ailleurs : la règle qui détecte une vérification de certificat désactivée
est \textbf{volontairement non écartée}, son constat étant neutralisé par des annotations
nominatives posées là où le comportement est assumé, de sorte que \textbf{toute nouvelle occurrence
ailleurs dans le dépôt fasse échouer la chaîne} --- une exclusion globale aurait supprimé le constat
existant et, du même geste, la capacité à détecter le suivant.

\begin{alertbox}[boxsep=1pt, left=4pt, right=4pt, top=3pt, bottom=3pt]
Une limite de cet épinglage est reconnue et mesurée : la version de l'analyseur est figée, mais le
\emph{jeu de règles} qu'il télécharge est distant et mutable. Deux exécutions à version d'outil
identique ont dénombré 612 règles le 20/08/2026 et 606 le 24/08/2026. Le verdict de la chaîne
dépend donc en partie d'un référentiel que le dépôt ne contrôle pas ; la clôture de cet écart
suppose de figer aussi le jeu de règles.
\end{alertbox}

\subsection{Authentification sans secret permanent}
\label{subsec:wif}

L'étape 1 réalise l'exigence EX7 par la fédération d'identité décrite au \cref{ch:conception} : la chaîne ne
détient aucune clé, et cette absence n'est pas une intention --- la recherche exhaustive des
ressources de clé de compte de service, des fichiers d'identifiants et de la variable
d'environnement qui les désigne ne relève aucune occurrence dans le dépôt. Le point à valoriser est
la forme de la condition d'attribut qui restreint cette fédération : elle porte \textbf{à la fois}
sur le dépôt d'origine \textbf{et} sur la référence de branche
(annexe~\ref{ann:extraits-realisation}). Sans la contrainte de branche, n'importe quelle branche du
dépôt obtenait l'identité de déploiement et, avec elle, le droit de déployer une révision arbitraire
des services. Le garde-fou qui semble couvrir ce cas, la condition de branche posée sur l'étape de
construction, se trouve \textbf{côté forge} : déclaré dans le même fichier que ce qu'il protège, un
contributeur capable de modifier ce fichier le lève en même temps qu'il l'exploite. \textbf{Seul le
filtre du fournisseur cloud fait autorité}, parce qu'il est déclaré ailleurs, appliqué par une autre
partie et non modifiable depuis le dépôt --- application directe du principe d'identité (PR2) : la
décision d'autorisation appartient au vérificateur, jamais à l'émetteur.

Trois réserves complètent ce constat sans l'annuler : la liaison d'usurpation est posée sur
l'attribut de dépôt et non sur celui de référence, de sorte que la restriction de branche ne tient
que par la condition d'attribut ; la référence autorisée a une valeur par défaut commune aux deux
environnements ; et les trois paliers d'identité prévus à l'origine ont été remplacés par une
fédération et une identité de déploiement uniques par projet --- le chemin d'autorisation y gagne en
lisibilité, sans que l'impossibilité d'un accès croisé soit prouvée.

\subsection{Déploiement, retour arrière et frontière avec l'infrastructure}

Le déploiement de l'étape 7 remplace intégralement le trafic du service par la nouvelle révision,
sans répartition progressive ; le bloc de répartition a été rendu explicite dans le code et placé
hors du champ comparé par la chaîne d'infrastructure pour préparer un déploiement graduel, capacité
non exploitée à ce jour. Le retour arrière est donc un geste manuel documenté --- replacer la
totalité du trafic sur la révision précédente, restée disponible --- dont la condition de bascule
est la première mise en production, où il devient un délai de rétablissement non maîtrisé. La
vérification de disponibilité qui suit interroge le service par son domaine public, donc à travers
le filtrage applicatif : elle éprouve la chaîne d'accès complète, non le seul conteneur.

\begin{alertbox}[boxsep=1pt, left=4pt, right=4pt, top=3pt, bottom=3pt]
\textbf{Un défaut d'ordonnancement est rapporté ici plutôt que tu.} L'image est publiée au
registre \emph{avant} l'analyse de vulnérabilité : une image porteuse d'une faille critique
corrigeable entre donc dans le registre de confiance, seul son déploiement étant refusé. La
chaîne écrite plus tard pour l'application hébergée construit localement, analyse, puis ne
publie que si l'analyse passe --- c'est l'ordre correct, et il sert de cible d'alignement.
\end{alertbox}

Une réserve reste ouverte sur le schéma de données : la chaîne du socle n'automatise aucune
migration, alors que celle de l'application hébergée l'exécute comme une tâche distincte, attendue
jusqu'à son terme avant tout déploiement, et interdit à chaque instance de la déclencher elle-même
--- le bon motif, puisque $N$ instances qui migrent sans verrou de schéma produisent une corruption
silencieuse.

\subsection{Boucle entre livraison et détection : la huitième étape}
\label{subsec:boucle-f6}

L'étape 8 réalise le flux F6, hypothèse d'apport exposée au \cref{ch:conception} et rattachée à la
question Q6 du \cref{ch:cadre-existant} comme au verrou V4 du \cref{ch:etat-art} : elle exporte le résultat de l'analyse
de vulnérabilité vers une table dédiée, avec un droit d'écriture limité à cette seule table, de
sorte que gravité théorique et techniques observées coexistent enfin dans un entrepôt où elles
peuvent se parler. La chaîne est linéaire --- analyse de vulnérabilité, service d'encodage, table
des vulnérabilités, rapprochement avec les détections R1--R7 par technique ATT\&CK, ordre de
correction --- et son apport se dit sans emphase : \textbf{F6 ne couvre aucune exigence de
sécurité, il produit un ordre de traitement et non une réduction de risque} ; sa chaîne de
données existe, son utilité reste non validée (\cref{sec:eval-enrichissement}) tant que l'ordre
produit n'est pas comparé à l'ordre par gravité seule. Une limite s'y ajoute : le chargement
inscrit une technique nulle pour chaque vulnérabilité, si bien que le croisement fonctionne sur la
chaîne de données et non sur les valeurs. La condition de clôture est le renseignement de cette
colonne à la source.

\subsection{Refus vérifiés}
\label{subsec:refus}

Une chaîne de livraison doit être éprouvée par des cas de refus, faute de quoi rien ne distingue une
porte qui laisse passer de bons artefacts d'une porte qui ne fonctionne pas --- l'incident d'analyse
statique de la \cref{subsec:incidents} montre que la distinction n'est pas théorique. Trois
protocoles définissent ces cas : un secret factice doit arrêter la chaîne à la deuxième étape (T10),
une image de base porteuse d'une faille critique connue à la cinquième (T8), et un espace de
stockage volontairement exposé doit être refusé avant toute création (\cref{sec:chaine-infra}). La
campagne du \cref{ch:validation} les a départagés (\cref{sec:campagne}) : \textbf{T8 est conforme}, sur un
refus qui n'a pas été provoqué --- la porte a réellement arrêté une livraison sur une vulnérabilité
critique --- tandis que \textbf{T10 reste partiellement conforme}, la porte étant bloquante à chaque
livraison et sur l'historique complet, mais aucun refus n'ayant été opposé. Un quatrième cas, le
défaut de code détecté par analyse statique, est établi autrement et mieux : par la remise en
service de la porte le 19/08/2026 et sa première exécution réelle (\cref{subsec:incidents}).

Un cinquième refus, non planifié, est \textbf{réellement survenu sur le code écrit pour ce projet} :
la porte d'analyse de vulnérabilité a bloqué une livraison sur une faille critique embarquée par le
\emph{gestionnaire de paquets} de l'image de base, composant dont l'artefact final n'a aucun usage à
l'exécution (T8, \cref{sec:preuve-t11}).
Entre inscrire la vulnérabilité sur une liste d'exception et retirer la cause, c'est la seconde voie
qui a été retenue, la surface d'attaque disparaissant avec l'outil qui la portait : une porte a
trouvé, dans le socle lui-même, un défaut que sa conception n'avait pas anticipé.

\section{Chaîne de provisionnement de l'infrastructure}
\label{sec:chaine-infra}

La seconde chaîne obéit à une logique différente, et cette différence est une décision
d'architecture plutôt qu'un degré d'achèvement : elle est \textbf{automatisée jusqu'au contrôle et
humaine au-delà}. Son intitulé l'annonce --- le fichier qui la décrit s'appelle « contrôles
d'infrastructure », non « déploiement d'infrastructure ». Le déclenchement est filtré sur le
répertoire d'infrastructure, selon la même règle que la chaîne applicative, puis trois étapes
s'exécutent (annexe~\ref{ann:extraits-realisation}, \cref{fig:chaine-infrastructure}).

\begin{figure}[htbp]
\centering
\begin{tikzpicture}[
  x=1cm, y=1cm, font=\footnotesize,
  et/.style={nsysteme, text width=4.30cm, inner sep=3pt, minimum height=0.80cm},
  arret/.style={nabsent, text width=2.20cm, inner sep=2.5pt,
                minimum height=0.52cm, font=\scriptsize}]

% ============ frontiere automatique / manuel ================================
\draw[frontiereplan] (7.55,4.10) -- (7.55,9.75);
\node[font=\footnotesize\bfseries, anchor=south] at (3.10,9.95) {automatique};
\node[font=\footnotesize\bfseries, anchor=south, align=center, text width=5.4cm]
      at (11.60,9.95) {manuel --- identifiants de l'administrateur};

% ============ partie automatique ============================================
\node[et] (dep) at (3.10,9.10) {\textbf{1.} dépôt};
\node[et] (fmt) at (3.10,8.10) {\textbf{2.} format (récursif, bloquant)};
\node[et] (val) at (3.10,6.90) {\textbf{3.} validation --- deux environnements, sans accès distant};
\node[ncontrole, text width=4.30cm, inner sep=3pt, minimum height=0.80cm]
      (ana) at (3.10,5.15) {\textbf{4.} analyse de configuration --- \textbf{bloquante} sur gravité critique};
\draw[fluxdonnee] (dep) -- (fmt);
\draw[fluxdonnee] (fmt) -- (val);
\draw[fluxdonnee] (val) -- (ana);
\node[arret] (ar) at (3.10,3.95) {chaîne arrêtée};
\draw[fluxrefuse] (ana.south) -- node[croixrefus]{} (ar.north);

% ============ partie manuelle ===============================================
\node[nporte, text width=2.55cm] (por) at (11.60,5.15) {\textbf{5.} porte : revue du plan};
\node[et, text width=3.60cm] (app) at (11.60,7.20) {\textbf{6.} application};
\node[et, text width=3.60cm] (der) at (11.60,8.60) {\textbf{7.} détection de dérive planifiée};
\draw[fluxdonnee] (ana.east) -- (por.west);
\draw[fluxcontrole] (por.north) -- (app.south);
\draw[fluxdonnee] (app.north) -- (der.south);

% ============ retour de derive vers le depot ================================
\draw[fluxevenement] (der.north) -- (11.60,9.75)
      -- node[etiqsec, pos=0.55] {dérive constatée} (3.10,9.75) -- (dep.north);

% ============ legende interne ===============================================
\node[legendecadre, text width=15.00cm, anchor=north west, font=\scriptsize]
     at (0.05,3.35)
  {\legendetrait{fluxdonnee}{étape suivante}\quad
   \legendetrait{fluxcontrole}{décision de politique : le plan n'est appliqué qu'après revue}\quad
   \legendetrait{fluxevenement}{retour asynchrone}\\[2pt]
   \legendecroix{refus : la chaîne s'arrête}\quad losange et cadre épais : porte bloquante\quad
   trait vertical épais : frontière de privilège};

\end{tikzpicture}
\caption[Chaîne de provisionnement de l'infrastructure]{Chaîne de provisionnement de
l'infrastructure : la frontière automatique/manuel est un \emph{choix de privilège}, non un
degré d'achèvement --- la partie automatisée n'a, par construction, aucun accès au
fournisseur cloud.}
\label{fig:chaine-infrastructure}
\end{figure}

\textbf{La vérification de format est récursive et bloquante.} Rendre bloquant un contrôle de forme
peut surprendre ; c'est la condition d'une revue de plan utile, car une différence de mise en forme
et une différence de sens se présentent à l'identique dans un écart calculé.

\textbf{La validation syntaxique s'exécute sur les deux environnements, en initialisation sans
accès distant} : elle ne contacte ni le stockage de l'état ni le fournisseur cloud, et valide la
cohérence interne de la description --- références, types, arguments obligatoires, versions de
fournisseurs --- sans aucune autorisation sur l'infrastructure réelle. C'est ce qui permet à ce flux
de ne déclarer aucune permission d'émission de jeton : il ne peut structurellement pas
s'authentifier auprès du fournisseur cloud, et cette impossibilité est une propriété, pas un manque.

\textbf{L'analyse de sécurité de la description est une porte bloquante sur gravité critique}, sur
des motifs qui recoupent les trois interdictions que les politiques d'organisation auraient dû
porter (\cref{sec:env-travail}) ; un second passage, informatif, remonte les gravités inférieures
sans arrêter la chaîne, de sorte que la porte n'accumule pas de dette silencieuse. \textbf{Aucun
composant n'a été ajouté pour cette analyse} : c'est l'outil qui analyse les images, employé dans un
autre mode --- application du principe de justification (PR1).

\subsection{La revue du plan comme porte humaine}

Le plan de modification est produit puis appliqué par l'administrateur, hors de la chaîne, avec ses
propres identifiants. Automatiser l'application supposerait d'accorder à l'identité de déploiement
des droits de \emph{lecture et d'écriture étendus sur tout le projet} --- un plan complet lit toutes
les ressources gérées, une application complète les modifie --- soit la concentration de privilège
que le \cref{ch:conception} refuse.

\begin{keybox}[boxsep=1pt, left=4pt, right=4pt, top=3pt, bottom=3pt]
\textbf{Décision retenue.} La revue humaine du plan est la porte de cette chaîne.
\textbf{Condition de bascule :} l'automatisation devient rentable lorsque le périmètre justifie
une \emph{identité de déploiement d'infrastructure séparée} de l'identité de livraison
applicative --- c'est-à-dire à partir du troisième environnement, où la fréquence des
modifications et le nombre d'opérateurs rendent le geste manuel plus risqué que le privilège
qu'il évite.
\end{keybox}

Deux principes gouvernent l'ensemble. \textbf{Le plan est revu comme du code}, et se lit
intégralement, jamais par sa fin : l'incident rapporté en \cref{subsec:incidents} montre qu'un plan
à périmètre restreint peut contenir un remplacement destructif sans rapport.
\textbf{L'application est la seule voie d'écriture} : toute modification venue d'un autre chemin
devient une dérive, détectée à la planification suivante (\cref{subsec:derive}). D'où le verdict de
T18, partiellement conforme au \cref{ch:validation} : la revue a lieu, l'approbation n'est pas outillée.

\subsection{Protection contre les destructions}

Une chaîne de provisionnement peut détruire aussi vite qu'elle construit. Deux ressources seulement
portent le verrou de destruction du langage de description : les deux clés de chiffrement. Les
services d'exécution, les tâches et les modules applicatifs portent explicitement le verrou ouvert,
cohérent avec leur caractère reconstructible. L'instance de base de données et les tables de
l'entrepôt --- \emph{y compris celles qui portent les preuves} --- ne portent en revanche pas la
protection de suppression du fournisseur, le seul garde-fou de l'ensemble de données étant son refus
de disparaître tant qu'il contient des tables.

\begin{alertbox}[boxsep=1pt, left=4pt, right=4pt, top=3pt, bottom=3pt]
Cet écart est classé au premier rang du plan de correction pour un motif de proportion : le coût est
de quelques minutes, et l'effacement de la table des détections par une commande d'administration
les 24 et 25/08/2026 a établi que le risque n'est pas théorique. La correction pose la protection de
suppression du fournisseur sur l'instance de base de données et sur les trois tables de preuve ---
détections, verdicts d'analyste, événements de sécurité --- et sert le principe d'intégrité de la
preuve (PR3), qui perd son sens si la preuve peut être détruite d'un geste.
\end{alertbox}

Le quatrième scénario de refus porte sur une ressource volontairement mal configurée, dont le
blocage doit être constaté avant toute création plutôt que lors d'un audit a posteriori. Sa
définition est complète --- ressource, politique attendue, critère d'acceptation --- mais elle ne
porte pas de protocole numéroté : à la différence de T8 et de T10, elle n'entre pas dans la campagne
du \cref{ch:validation}, et son premier passage est porté par le plan de correction.

\subsection{Conduite d'une migration sur un système en service}
\label{subsec:migration-reseau}

La refonte réseau arrêtée sous la décision D15 (\cref{subsec:reseau-cible}) a été appliquée par
cette chaîne --- un plan lu intégralement puis appliqué à la main à chaque étape --- en douze étapes,
entre le 11 et le 25/08/2026. Elle est rapportée ici parce qu'elle est la seule opération du projet
à avoir déplacé le chemin du trafic d'un système déjà en service : l'ordre y compte davantage que le
contenu.

\textbf{Cet ordre obéit à une règle unique} --- \emph{on n'interdit jamais avant d'avoir observé,
et on n'observe qu'après avoir fait passer le trafic par l'endroit où on l'observe}. Les trois
sous-réseaux par locataire ont donc été créés d'abord (11/08) : ajout pur, qu'aucune ressource ne
référençait encore. Les charges de travail ont ensuite basculé une par une en sortie réseau
directe (12 au 14/08), en commençant par la tâche planifiée d'enrichissement --- dont l'échec eût
été visible et sans effet sur un utilisateur --- et en gardant pour la fin le service le plus
exposé ; chaque bascule produit une révision, et une révision qui ne démarre pas laisse la
précédente servir, de sorte que le retour arrière est un redéploiement et non une reconstruction.
Le connecteur d'accès sans serveur n'a été supprimé qu'ensuite (15/08), une fois qu'aucune
ressource ne le référençait et que le plan ne présentait qu'une seule destruction.

\textbf{Sept jours d'observation ont remplacé le pari par la mesure} (15 au 22/08) : échantillonnage
des journaux de flux porté à un paquet sur deux, journalisation de la passerelle de traduction
d'adresses portée à toutes ses connexions, et construction depuis l'entrepôt de l'inventaire exhaustif
des destinations réellement jointes. C'est cet inventaire, et non une hypothèse sur les besoins des
charges, qui a dicté les six autorisations posées le 22/08 ; c'est la même fenêtre qui a établi que
la passerelle ne traitait aucun trafic, et justifié sa suppression par une mesure.

\textbf{Une seule étape était réellement dangereuse, et elle a été traitée comme telle} : la pose
du refus par défaut en sortie, le 23/08, qui coupe tout ce que les cinq autorisations n'ont pas
nommé. Elle a été jouée hors heures ouvrées, la commande de retour arrière préparée avant
l'application, et n'a été déclarée réussie qu'au terme de trente minutes d'observation. Le risque
résiduel avait été nommé avant d'être couru : un flux légitime qui ne survient pas pendant la
fenêtre de sept jours --- une tâche mensuelle, un chemin d'erreur --- n'entre dans aucun
inventaire. C'est la raison d'être des refus nommés posés juste avant lui : ils font apparaître un
tel flux dans les journaux sous un nom de règle lisible, au lieu de le laisser disparaître dans un
refus général qui n'explique rien. La passerelle, son routeur, les deux anciens sous-réseaux vides
et les trois règles d'entrée qui ciblaient une étiquette que rien ne portait ont été retirés le
lendemain ; le régime permanent --- échantillonnage ramené à un paquet sur dix, métadonnées
exclues --- a été posé le \dateref.

\section{Chaîne de détection}
\label{sec:chaine-detection}

La collecte (flux F4) rassemble vers l'entrepôt de supervision les journaux d'audit du fournisseur,
ceux des services d'exécution et les verdicts du filtrage applicatif, avec un filtrage appliqué
\textbf{à la source} et non après stockage (BNF3). Une réserve porte sur les journaux d'audit : ceux
d'activité d'administration sont collectés avec les autres sources, tandis que ceux d'accès aux
données --- entrepôt, secrets, clés, base applicative --- restent dans l'espace par défaut du
fournisseur, en conservation de trente jours, modifiable. L'attribution d'une action privilégiée
visée par EX13 n'est donc établie que pour le plan d'administration, d'où l'exécution du protocole
T13 sur les deux plans.

L'entrepôt compte dix tables (\cref{tab:modele-donnees}), dont six voient leur identité d'écriture
restreinte à une table unique par une liaison de droits : cette séparation vise l'exigence EX11 sans
rendre les preuves immuables --- une identité d'administration ou une liaison IAM excessive peut
encore les modifier, ce qui doit être contrôlé séparément. Le score d'incident, recalculé à la
lecture et jamais persisté (\cref{ch:conception}), est une \textbf{heuristique de tri} et non un verdict : il
additionne un poids par gravité et une prime fixe lorsque les détections d'une même entité relèvent
d'au moins deux tactiques distinctes, signe d'une chaîne d'attaque probable. L'asymétrie est
voulue : le score est recalculé à chaque consultation et jamais conservé, le verdict de l'analyste
est conservé et jamais modifié --- traduction du principe d'intégrité de la preuve (PR3),
\textbf{ce que la machine ordonne est révisable, ce que l'humain juge est définitif} ; la colonne
« verdict » reste vide tant que l'analyste n'a pas tranché.


Trois colonnes portent des données à caractère personnel : l'adresse source et l'identifiant
d'utilisateur des journaux d'accès, et l'entité des détections, le plus souvent une adresse. Les
journaux expirent à quatre-vingt-dix jours ; les détections, qui portent la même adresse, n'expirent
pas --- écart assumé au titre de la valeur probante, à trancher par le responsable de traitement.
Les 872 vecteurs de référence couvrent 697 techniques et sous-techniques du référentiel ATT\&CK
Entreprise \cite{mitre_attack}. Chaque locataire dispose en outre d'un espace de stockage objet
dédié, chiffré par clé gérée, protégé contre tout accès public et administrable en écriture par la
seule identité de l'application concernée : il porte les \textbf{enregistrements vocaux} de
l'application pilote, données biométriques qui fondent la contrainte réglementaire du \cref{ch:cadre-existant}.

\begin{alertbox}[boxsep=1pt, left=4pt, right=4pt, top=3pt, bottom=3pt]
La journalisation des refus du pare-feu a été activée et vérifiée le 19/08/2026, ce qui clôt
l'écart É1 et rend la source de signal disponible ; elle porte, depuis la migration rapportée en
\cref{subsec:migration-reseau}, sur les six règles de refus du réseau --- dont les quatre refus
croisés entre locataires --- que le puits existant collecte sans modification. Le trajet complet d'un refus jusqu'à l'entrepôt a été
mesuré depuis : le 24/08/2026, trois sorties refusées --- externe, croisée entre locataires et vers
la base --- ont été retrouvées dans l'entrepôt en 2~min~55~s à 4~min~12~s, ce qui rend le
protocole T20 \textbf{conforme} (\cref{sec:campagne}) ; sa cadence de rejeu est trimestrielle.
\end{alertbox}

Il n'existe pas de convertisseur automatique depuis le format neutre de règles de détection vers
le dialecte de l'entrepôt retenu (\cref{ch:etat-art}). La traduction des sept règles R1--R7 a donc été
réalisée manuellement, selon un gabarit systématique --- source, condition de sélection, fenêtre
temporelle, agrégation, identifiants, faux positifs connus --- dont
l'annexe~\ref{ann:extraits-realisation} donne un exemple.

La chaîne complète, de la collecte à l'affichage, est celle conçue au \cref{ch:conception}
(\cref{fig:cycle-donnee}) et la réalisation n'en a pas déplacé une étape : les journaux bruts
alimentent les sept requêtes planifiées R1--R7 sur une fenêtre de 15~min, les détections
qu'elles écrivent sont lues par l'API qui calcule le score à la requête, et la tâche
d'enrichissement forme un \textbf{second cycle indépendant} --- encodage puis recherche
vectorielle à trois candidats --- dont le résultat n'est joint à aucun score. Les deux cycles
étant indépendants, une détection est affichée dès qu'elle est écrite et son enrichissement
arrive au cycle suivant.

Le plafond que cette traduction manuelle impose à la couverture est établi en
\cref{sec:detection-code} (verrou V2) et mesuré au \cref{ch:validation} ; la contrepartie est que les règles
sont du code, versionnées, revues et déployées par la chaîne d'infrastructure.

\section{Couche d'enrichissement sémantique}

Les contrôles appliqués aux poids du modèle de représentation de phrases retenu
\cite{securebert2022,smet2023}, traités comme un artefact tiers non fiable, sont ceux du
\cref{ch:conception}. Ce que la réalisation ajoute est une mesure : \textbf{la fidélité est vérifiée, pas
supposée}. La quantification en entiers, testée puis rejetée par la porte d'acceptation codée,
dégradait les représentations au point de ne produire \textbf{aucune correspondance correcte de
premier rang sur cinq cas} ; le modèle livré est donc en précision flottante simple. La tâche
d'enrichissement retient trois candidats plutôt qu'un, afin que l'analyste compare des alternatives
au lieu de recevoir une affirmation unique.

\textbf{Posture de sortie réseau obtenue.} Depuis la migration rapportée en
\cref{subsec:migration-reseau}, les charges portent leurs interfaces de sortie \emph{dans} un
sous-réseau : tout leur trafic sortant traverse le pare-feu, le refus par défaut en sortie est posé,
et cinq autorisations seulement le percent, chacune ciblée par une identité de service et par un
port. Le plan d'adressage, les douze règles et les trois lectures qu'elles appellent sont établis en
\cref{subsec:reseau-cible} ; ce que la réalisation ajoute est que \textbf{le pare-feu est devenu une
source de journal nommée} --- les refus croisés entre locataires et l'interdiction faite au tableau
de bord de joindre la base portent chacun un nom de règle qui dit ce qui a été tenté, et rejoignent
l'entrepôt par le puits déjà en place, \textbf{sans qu'aucune modification de la chaîne de détection
ait été nécessaire} (annexe~\ref{ann:extraits-realisation}).

Cette posture n'établit pas l'impossibilité d'une exfiltration : elle en ferme la voie directe et en
laisse deux ouvertes, énoncées en \cref{subsec:reseau-cible} --- les interfaces du fournisseur, et
les deux charges détachées du réseau. La seconde n'est pas une hypothèse : le 25/08/2026, une sonde
empruntant le chemin de la révision réellement déployée du service d'encodage a atteint une
destination externe en 412~ms, quand la même sonde passant par le sous-réseau était refusée deux
fois. \textbf{Le protocole T14 est donc non conforme} ; le correctif --- réintégrer les deux charges
au sous-réseau sous une étiquette qu'aucune autorisation de sortie ne couvre --- coûte une révision
sans changement applicatif. Le trajet d'un refus jusqu'à l'entrepôt est, lui, mesuré et
\textbf{conforme} par T20 (24/08/2026, \cref{sec:campagne}).

\subsection{État d'avancement réel de la couche}
\label{subsec:avancement-enrichissement}

L'écart entre conception et réalisation est ici le plus important
(\cref{tab:avancement-enrichissement}).

\begin{table}[htbp]
\caption[État d'avancement de la couche d'enrichissement]{État d'avancement de la couche
d'enrichissement sémantique}
\label{tab:avancement-enrichissement}
\centering
\scriptsize
\setlength{\tabcolsep}{4pt}
\begin{tabular}{|P{5.20cm}|P{7.35cm}|P{2.55cm}|}
\hline
\textbf{Composant} & \textbf{Réalisé} & \textbf{Statut} \\
\hline
Préparation et figement du modèle & Oui & Vérifié \\ \hline
Chemin public direct absent dans la configuration & Oui --- mais la sortie \emph{managée} de la
révision déployée reste hors du pare-feu & \textbf{T14 non conforme} (25/08/2026) \\ \hline
Table de vecteurs de référence et index & Oui & Vérifié \\ \hline
Tâche d'enrichissement, recherche vectorielle & Oui & Réalisé \\ \hline
Cadence prévue (5~min) & Non --- 15~min & Écart assumé (coût) \\ \hline
Cache par étiquette SHA d'entrée & Calculé, jamais exploité & Non réalisé \\ \hline
Vue de corrélation entre alertes & Non & Non réalisé \\ \hline
Enrichissement dans le score d'incident & Non & Écart majeur \\
\hline
\end{tabular}
\end{table}

\begin{alertbox}[boxsep=1pt, left=4pt, right=4pt, top=3pt, bottom=3pt]
\textbf{La couche d'enrichissement produit ses résultats, mais ceux-ci n'alimentent pas le score
d'incident présenté à l'analyste.} Cet écart ne remet pas en cause la démonstration technique --- la
chaîne complète fonctionne et le \cref{ch:validation} la mesure --- mais il borne l'affirmation possible : le
rattachement automatique produit des candidats consultables sans preuve suffisante de précision, de
coût stabilisé ni d'amélioration de la priorisation.
\textbf{Décision du \dateref{} :} le positionnement informationnel est assumé comme état livré, et
la jointure au score est reportée en perspective. Elle ne se réduit pas à une jointure technique :
elle suppose une règle de décision évaluée et surtout l'ajout d'une clé d'entité à la table
d'enrichissement, dont le \cref{ch:conception} montre qu'elle est absente du schéma.
\end{alertbox}

\section{Tableau de bord et API de supervision}
\label{sec:supervision}

La supervision repose sur deux logiciels écrits pour ce projet : le tableau de bord, application web
à rendu serveur, et l'API de supervision qui l'alimente (\cref{fig:composants-logiciel}). Tout passe
par l'API, qui applique le contrôle d'autorisation ; cette séparation réalise les exigences EX5 et
EX6 et évite l'erreur fréquente donnant à une interface web les droits de lecture d'un entrepôt. Les
vues se rafraîchissent en continu et signalent la fraîcheur de leurs données, mais cette
actualisation porte sur des résultats \emph{déjà calculés} : l'interface est vive, la détection ne
l'est pas au même degré.

\begin{figure}[htbp]
\centering
\begin{tikzpicture}[
  x=1cm, y=1cm, font=\footnotesize,
  comp/.style={nsysteme, text width=3.15cm, inner sep=3pt, minimum height=1.15cm},
  cle/.style={ncontrole, text width=3.15cm, inner sep=3pt, minimum height=1.15cm,
              fill=menalplanB},
  ext/.style={nstock, text width=2.25cm, inner sep=3pt, minimum height=1.15cm},
  titre/.style={font=\footnotesize\bfseries, align=center, text width=4.2cm},
  cadre/.style={draw=menalencre, line width=0.6pt, densely dashed, rounded corners=4pt}]

% ============ tableau de bord (conteneur de gauche) =========================
\node[cle]  (mw)    at (4.10,7.10) {Garde de session\\\gcprod{intergiciel}};
\node[comp] (pages) at (4.10,5.45) {12 vues, dont 10 rendues au serveur};
\node[cle]  (bff)   at (4.10,3.75) {7 routes serveur : jeton déposé en \emph{cookie} hors de portée du script};
\node[comp] (cli)   at (4.10,1.85) {Client d'appel typé et 15 composants d'affichage};

% ============ API de supervision (conteneur de droite) ======================
\node[cle]  (mwapi) at (9.30,7.10) {En-têtes de sécurité et journal d'audit};
\node[cle]  (auth)  at (9.30,5.45) {Authentification : jeton signé, second facteur, rôle};
\node[comp] (siem)  at (9.30,3.75) {Routeur de supervision\\\textbf{9 points d'entrée}};
\node[comp] (crud)  at (9.30,1.85) {Comptes, journaux, alertes\\\textbf{4 points d'entrée}};

% ============ acteur et plan de donnees =====================================
\node[nacteur, text width=1.85cm, minimum height=1.05cm] (analyst) at (0.95,5.45) {Analyste\\(navigateur)};
\node[ext] (sm)  at (13.55,7.10) {Coffre de secrets};
\node[ext] (bq)  at (13.55,3.75) {Entrepôt de supervision};
\node[ext] (sql) at (13.55,1.85) {Base applicative};

% ============ appels synchrones =============================================
\draw[fluxdonnee] (analyst.north) |- (mw.west);
\draw[fluxcontrole] (mw) -- node[etiqsec] {session} (pages);
\draw[fluxdonnee] (pages) -- (bff);
\draw[fluxdonnee] (bff)   -- (cli);
\draw[fluxdonnee] (mwapi) -- (auth);
\draw[fluxcontrole] (auth) -- node[etiqsec] {rôle vérifié} (siem);
\draw[fluxdonnee] (siem)  -- (crud);
\draw[fluxdonnee] (bff.east) -- ++(0.28,0) |- node[etiq, pos=0.72] {identifiants} (auth.west);
\draw[fluxdonnee] (cli.east) -- ++(0.60,0) |- node[etiq, pos=0.72] {jeton porteur} (siem.west);
\draw[fluxdonnee] (auth.east) -- ++(0.35,0) |- (sm.west);
\draw[fluxdonnee] (siem.east) -- node[etiq] {lecture seule} (bq.west);
\draw[fluxdonnee] (crud.east) -- (sql.west);

% ============ ecriture au mieux, hors du chemin de reponse ==================
\draw[fluxevenement] (mwapi.east) -- (15.30,7.10) -- node[etiqsec, pos=0.55] {audit}
      (15.30,1.85) -- (sql.east);

% ============ frontieres de deploiement et plan de donnees ==================
\begin{scope}[on background layer]
  \node[cadre, fit=(mw)(pages)(bff)(cli), inner sep=6pt] (fd) {};
  \node[cadre, fit=(mwapi)(auth)(siem)(crud), inner sep=6pt] (fa) {};
  \node[bandedonnee, fit=(sm)(bq)(sql), inner sep=5pt] (fs) {};
\end{scope}
\node[titre, anchor=south] at (fd.north) {Tableau de bord\\\gcprod{3\,408 lignes TypeScript}};
\node[titre, anchor=south] at (fa.north) {API de supervision\\\gcprod{1\,855 lignes Python}};
\node[titrebande, anchor=south] at (fs.north) {plan de données};

% ============ legende interne ===============================================
\node[legendecadre, text width=15.00cm, anchor=north west, font=\scriptsize]
     at (-0.15,0.55)
  {\legendetrait{fluxdonnee}{appel synchrone}\hfill
   \legendetrait{fluxcontrole}{point d'application du contrôle d'accès}\hfill
   \legendetrait{fluxevenement}{écriture au mieux, hors du chemin de réponse}\\[2pt]
   cadre épais, fond grisé : composant portant un contrôle\hfill
   cadre en tirets : un conteneur\hfill
   double filet : dépôt\hfill
   bande claire : plan de données};

\end{tikzpicture}
\caption[Architecture interne des deux logiciels écrits]{Architecture interne des deux logiciels
écrits : composants, frontières de déploiement et points d'application du contrôle d'accès. Le rôle
est vérifié sur les treize points d'entrée de données de l'API --- neuf de supervision et quatre de
gestion --- et jamais dans le tableau de bord, qui ne détient aucune autorisation propre ; l'échec
de l'écriture d'audit est absorbé et ne fait jamais échouer la requête de l'analyste.}
\label{fig:composants-logiciel}
\end{figure}

\subsection{La chaîne de données d'une vue}
\label{subsec:chaine-donnees-vue}

Le point le plus solide de cette architecture n'est pas visible à l'écran : \textbf{aucune vue du
tableau de bord n'interroge directement l'entrepôt de supervision ni la base applicative}. Toutes
empruntent les cinq mêmes sauts, sans exception vérifiée ligne à ligne
(\cref{fig:chaine-donnees-vue}, extrait en annexe~\ref{ann:extraits-realisation}) : cette uniformité
est ce qui rend la chaîne auditable.

\begin{figure}[htbp]
\centering
\begin{tikzpicture}[
  x=1cm, y=1cm, font=\footnotesize,
  saut/.style={nsysteme, inner sep=3pt, minimum height=0.85cm}]

% ============ rangee haute : navigateur, tableau de bord, bord ==============
\node[saut, text width=2.05cm] (nav) at (1.35,2.95) {\textbf{1.} navigateur};
\node[saut, text width=3.05cm] (tdb) at (7.35,2.95) {\textbf{2.} tableau de bord (composant rendu au serveur)};
\node[saut, text width=3.25cm] (brd) at (13.05,2.95) {\textbf{3.} bord : répartiteur et filtrage applicatif};
\draw[fluxdonnee] (nav) -- node[etiq, text width=2.15cm, align=center]
      {témoin inaccessible au script} (tdb);
\draw[fluxdonnee] (tdb) -- node[etiq, text width=2.05cm, align=center]
      {en-tête d'autorisation} (brd);

% ============ la frontiere que le jeton ne franchit pas =====================
\draw[seuil] (3.05,1.95) -- (3.05,4.10);
\node[font=\footnotesize\bfseries, anchor=south, align=center, text width=4.6cm]
      at (3.05,4.15) {le jeton ne franchit jamais cette frontière};

% ============ rangee basse : API puis entrepot ==============================
\node[saut, text width=3.25cm] (api) at (12.55,0.55) {\textbf{4.} API de supervision --- vérification du jeton et du rôle};
\node[saut, text width=3.05cm] (ent) at (7.35,0.55) {\textbf{5.} entrepôt de supervision, base applicative};
\draw[fluxcontrole] (brd.south) -- node[etiq, text width=2.35cm, align=center]
      {seconde traversée du bord} (api.north);
\draw[fluxdonnee] (api) -- node[etiq, text width=1.95cm, align=center]
      {requête paramétrée} (ent);

% ============ legende interne ===============================================
\node[legendecadre, text width=15.00cm, anchor=north west, font=\scriptsize]
     at (0.05,-0.35)
  {\legendetrait{fluxdonnee}{saut de la chaîne}\hfill
   \legendetrait{fluxcontrole}{traversée d'un point de contrôle}\hfill
   trait vertical en tirets : frontière que le jeton ne franchit pas};

\end{tikzpicture}
\caption[Chaîne de données d'une vue du tableau de bord]{Chaîne de données d'une vue du
tableau de bord : les cinq sauts, et la frontière que le jeton ne franchit pas. L'autorité
d'autorisation est l'API de supervision, qui vérifie le jeton et le rôle ; l'identité du
tableau de bord ne détient aucune autorisation sur les données.}
\label{fig:chaine-donnees-vue}
\end{figure}

Aucun script exécuté dans le navigateur ne peut donc lire le jeton, ni l'exfiltrer, y compris en cas
d'injection de script réussie dans une vue : \textbf{le tableau de bord n'est pas un client
privilégié, c'est un client comme un autre, qui n'a de pouvoir que celui du jeton qu'on lui a
confié}. Les sept routes serveur ne dérogent pas à ce motif : ce sont des relais
d'authentification, qui lisent le témoin, en font un en-tête d'autorisation, transmettent et rendent
la réponse. L'une d'elles porte une décision : le relais de connexion \emph{préserve} le code de
refus pour dépassement de seuil au lieu de le réécrire en refus d'authentification, ce qui aurait
fait croire à l'analyste qu'il s'était trompé de mot de passe et privé le centre d'opérations du
signal d'attaque. Le garde de session vérifie la présence du témoin et sa non-expiration par simple
décodage de la date d'échéance, \textbf{sans vérification de signature} --- le tableau de bord ne
connaît pas la clé et ne doit pas la connaître ; ce contrôle est \emph{documenté dans le code comme
un contrôle d'ergonomie et non d'autorisation}, la seule source de vérité restant l'API.

\subsection{Cartographie des vues}

Le \cref{tab:vues-dashboard} met en regard, pour chaque vue, ce qu'elle apporte à
l'analyste avec le point d'entrée qu'elle appelle, et sa capacité à être filtrée par locataire.


La colonne de droite est presque toujours mal lue : la portée du filtre par locataire n'est pas un
défaut d'interface qu'un développement complémentaire viendrait combler, c'est une \textbf{propriété
du modèle de données}. Une vue ne peut être filtrée par
locataire que si la table qu'elle interroge porte une clé de locataire fiable ; or trois tables
seulement de l'entrepôt en portent une, et une seule de ces trois clés est renseignée sur la
totalité des lignes (\cref{tab:modele-donnees}). Les vues
d'ensemble, de vulnérabilités et de santé des règles restent donc globalement fusionnées quel que
soit le locataire sélectionné, ce que le sélecteur affiché à l'écran documente lui-même. La
condition de bascule est l'introduction d'une clé de locataire obligatoire à l'ingestion, préalable
de tout cloisonnement au stockage.

\begin{alertbox}[boxsep=1pt, left=4pt, right=4pt, top=3pt, bottom=3pt]
Le filtre par locataire est porté par un témoin déposé côté navigateur, non signé, et le jeton ne
porte aucune revendication de locataire. Tout analyste authentifié peut donc changer de locataire
ou retirer le filtre. Le mémoire l'énonce comme le code l'énonce : \textbf{c'est un raffinement
de lecture, pas une frontière de sécurité}. Le transformer en frontière suppose une revendication
de locataire dans le jeton et un filtre imposé côté serveur --- travail identifié, chiffré, et
subordonné au préalable de la clé de locataire.
\end{alertbox}

\subsection{Authentification, second facteur et contrôle d'accès}
\label{subsec:authent-dashboard}

Trois rôles sont définis dans le jeton émis par l'API : administrateur, analyste et service ---
cette dernière identité, non humaine, n'étant référencée par aucun contrôle à ce jour ; aucun rôle
n'est prévu pour les utilisateurs de l'application hébergée, qui n'accèdent jamais au tableau de
bord. L'écriture d'un verdict illustre l'exception prévue au principe d'intégrité de la preuve
(PR3) : elle est possible, mais uniquement dans une table qui lui est propre
(\cref{tab:modele-donnees}), jamais dans les tables de preuve, et réservée à l'administrateur ---
arbitrage plus restrictif que la répartition des tâches d'un centre d'opérations, cohérent avec la
règle « une seule identité écrit dans une table donnée », et dont la condition de réévaluation est
l'arrivée d'un second analyste. L'interface a été servie à un compte analyste le 23/08/2026 sur des
détections réelles de la recette, une alerte R2 y apparaissant rattachée à la technique T1498 et
attribuée au locataire concerné.

\begin{figure}[!htbp]
\centering
\IfFileExists{figures/K05.png}{%
  \setlength{\fboxrule}{0.4pt}\setlength{\fboxsep}{0pt}%
  \fcolorbox{black!30}{white}{\includegraphics[width=0.82\textwidth]{figures/K05.png}}%
}{%
  \begin{minipage}{0.82\textwidth}
  \begin{extraitconf}[breakable=false, fontupper=\footnotesize,
      title={Capture K05 : spécification de prise}]
  \textbf{Vue :} tableau de bord de supervision, page \emph{Détections}, session ouverte sous un
  compte portant le rôle analyste --- pas le rôle administrateur, qui est l'objet même de la
  preuve.\par\smallskip
  \textbf{Date :} 23/08/2026 si la vue est repeuplée depuis la table d'archive de récupération, la
  table vive ayant été vidée par une suppression d'administration les 24 et 25/08/2026 ; à défaut,
  date du rejeu. \emph{La date portée par la capture fait foi.}\par\smallskip
  \textbf{Contenu probant :} l'interface décrite servie à un rôle analyste sur des
  détections réelles de la recette --- au moins une ligne de la règle R2 rattachée à la technique
  T1498 et portant le locataire dans sa colonne de service, l'horodatage de la détection et
  l'indicateur de fraîcheur visibles dans le cadre, et le sélecteur de locataire affichant lui-même
  la portée partielle de son filtre.\par\smallskip
  \textbf{Masquage} (rectangle opaque, jamais de flou) \textbf{:} identifiant de projet cloud,
  courriel du compte connecté, adresses réelles.
  \end{extraitconf}
  \end{minipage}%
}
\caption[Tableau de bord servi à un compte analyste, page des détections]{Tableau de bord de
supervision servi à un compte de rôle analyste, page des détections. À y lire : une détection de la
règle R2 rattachée à la technique T1498 et attribuée au locataire, l'horodatage de la détection
porté par la capture elle-même, et le sélecteur de locataire qui énonce la portée partielle de son
filtre.}
\label{fig:capture-K05}
\end{figure}

\textbf{Le facteur de connaissance.} Une longueur minimale de \textbf{douze caractères} est imposée
par le schéma de validation d'entrée et non par un contrôle écrit à la main : la contrainte est
portée par le contrat de l'API et vaut pour tout appelant, sans exigence de complexité, d'expiration
ni d'historique, conformément aux recommandations actuelles. La vérification borne en outre la
longueur de l'entrée avant la fonction de hachage, faute de quoi un appel \emph{non authentifié}
portant un mot de passe très long provoquait une erreur interne au lieu d'un refus --- défaut de
disponibilité atteignable sans compte, corrigé.

\textbf{La session.} Deux propriétés de la vérification du jeton signé décrit au \cref{ch:conception} méritent
d'être relevées. L'algorithme de signature est \textbf{figé en constante} et n'est jamais lu depuis
l'en-tête du jeton présenté, ce qui immunise contre la confusion d'algorithme --- la classe
d'attaque où un jeton signé par un autre algorithme, voire par aucun, est accepté parce que le
vérificateur croit ce que le jeton déclare de lui-même. La validation exige par ailleurs la présence
effective des champs d'échéance et d'émission, un jeton dépourvu d'échéance ayant été auparavant
valable indéfiniment.

\textbf{Le second facteur.} La séparation du jeton intermédiaire et du jeton d'accès, conçue au
\cref{ch:conception}, est \textbf{vérifiée à chaque appel} : la dépendance qui identifie l'appelant refuse
tout jeton dont le type n'est pas celui d'un jeton d'accès. Le point d'entrée qui reçoit le code
temporel n'était par ailleurs couvert que par un limiteur applicatif, non par la limitation de
débit du périmètre ; il a été ajouté aux chemins d'authentification protégés le 19/08/2026 et
vérifié en conditions réelles le jour même : onze requêtes successives ont produit dix refus de
validation, puis un refus pour dépassement de seuil. C'est cette trace qui étaye le test T2.

\textbf{L'anti-force-brute repose sur deux couches, et le choix de la clé de comptage est un
raisonnement d'ingénieur plutôt qu'une configuration.} La couche principale est au bord : le
filtrage applicatif limite les appels aux chemins d'authentification par adresse source réelle, avec
bannissement temporaire au dépassement. La couche applicative ne peut pas utiliser cette clé, le
tableau de bord et l'API partageant une \textbf{adresse de sortie commune} : lorsque la requête est
relayée par le tableau de bord, la seule adresse que voit le limiteur applicatif est celle de cette
sortie, de sorte qu'une clé par adresse rendrait le compteur \emph{commun à tous les analystes}. La
conséquence est quantifiable : douze requêtes non authentifiées suffisaient à verrouiller le centre
d'opérations entier, en refusant la connexion à des analystes qui n'avaient rien fait. La clé porte
donc sur le \textbf{défi} --- le jeton intermédiaire remis à l'issue du mot de passe --- ce qui
isole chaque tentative dans son propre budget d'essais. Le code énonce lui-même la limite
résiduelle : rejouer l'authentification par mot de passe réobtient un jeton intermédiaire, donc un
budget, ce qui reporte la charge sur la couche de bord, laquelle voit l'adresse réelle du client. La
leçon dépasse ce cas : \textbf{un contrôle qui suppose qu'une adresse réseau désigne un utilisateur
produit des blocages collectifs dès qu'une agrégation existe entre lui et le client} --- la question
n'est pas « quel seuil ? » mais « que compte-t-on, et cet objet est-il propre à celui que l'on veut
ralentir ? ».

Le secret partagé du second facteur était auparavant stocké en clair. Un \textbf{chiffrement
authentifié applicatif} a été mis en place le 19/08/2026, avec migration non destructive tolérant
les valeurs antérieures et clé conservée au coffre de secrets ; la migration des valeurs déjà en
base n'étant pas documentée comme exécutée, l'écart n'est pas clos.

\subsection{Le logiciel écrit : structure, tests et dette assumée}
\label{subsec:logiciel-ecrit}

Les deux composants de cette interface sont les seuls éléments du socle écrits ligne à ligne plutôt
que déclarés (\cref{tab:logiciel-ecrit}). Le contrôle d'autorisation conditionne la portée de
l'exigence EX5 et appelle une description exacte : les neuf points d'entrée de supervision déclarent
tous le \emph{même} composant, qui valide le jeton, en extrait le rôle et refuse l'appel si ce rôle
n'est pas attendu --- huit admettent l'administrateur et l'analyste, le neuvième, celui qui
enregistre un verdict, n'admet que l'administrateur. Il n'y a donc pas neuf contrôles à auditer mais
un seul, déclaré neuf fois, ce qui \emph{déplace} le risque : un point d'entrée ajouté sans la
déclaration serait ouvert sans qu'aucun test existant ne s'en aperçoive. Le remède, un test
d'inventaire vérifiant que toute route déclarée porte son contrôle, est porté au plan de
durcissement de l'API.


Deux lignes de ce tableau sont défavorables et appellent leur cause. La couverture de 3,00~\% tient
à ce que les vues ne sont pas testées : seules des fonctions pures le sont, et deux des quatre suites
éprouvent du code recopié dans le fichier de test plutôt que le module lui-même --- défaut de méthode
plus que défaut de couverture. Le progrès réel est ailleurs : la chaîne \emph{exécute} désormais ces
tests, en porte bloquante avant la construction. L'absence de version sur les points d'entrée est
sans conséquence tant que le tableau de bord en est le seul client, et sa condition de réévaluation
est l'ouverture à un client tiers.

Deux comportements de repli sont visibles à l'exécution et doivent être nommés. L'intergiciel qui
écrit le journal d'audit \textbf{absorbe toute exception} : une panne d'écriture ne fait jamais
échouer la requête de l'analyste. L'arbitrage est délibéré --- aveugler le centre d'opérations
pendant que la base est en difficulté serait pire --- mais il entre en tension avec l'intégrité de la
chaîne de preuve, une défaillance d'écriture ne se signalant par rien ; le contrôle compensatoire est
identifié, \textbf{une alerte sur l'absence d'écriture}, déclenchée lorsque le compteur de lignes
d'audit tombe à zéro alors que le service reçoit du trafic. Le tableau de bord \textbf{bascule sur un
jeu de données de démonstration} lorsque l'API est injoignable : repli signalé par un bandeau
visible, désactivé sur le flux d'authentification --- aucune session ne peut donc être ouverte sans
l'API --- et autorisé hors d'une construction de production seulement. Comme il porte un catalogue de
règles \emph{fictif}, la règle s'énonce sans détour : aucune démonstration hors de l'image de
production, la condition de clôture étant l'alignement du catalogue de repli sur les sept règles
réelles.

\section{Accueil d'un locataire et frontière de responsabilité}
\label{sec:locataire}
\label{sec:second-locataire}

Le besoin~BF13 exigeait que le socle accueille une nouvelle application sans reconception. Le module
d'exécution des services applicatifs a donc été rendu paramétrique en cours de projet : le nom de
l'application est une variable, et chaque application reçoit sa propre identité, ses propres secrets
et son propre service. Le socle héberge la plateforme MENAL elle-même, API et tableau de bord, et à
côté d'elle l'application pilote ELSON du \cref{ch:cadre-existant} : c'est cette seconde instance, « le second
locataire », qui sert de preuve d'accueil.

\subsection{La procédure d'accueil et la portée de l'isolation}

La procédure a été reconstituée à partir du cas réellement déployé, et non écrite \emph{a priori} :
chaque étape correspond à un geste qui a dû être fait. Elle suppose de l'application quatre prérequis
non négociables --- aucun état sur le disque local, configuration par variables d'environnement, un
secret par usage, écoute sur le port fourni avec un point de santé sans dépendance externe --- et se
déroule ensuite en six étapes (\cref{fig:accueil-locataire}). Quatre d'entre elles portent une
contrainte que le cas réel a imposée. L'\textbf{enregistrement de nom de domaine} doit exister
\emph{avant} tout provisionnement, faute de quoi le certificat géré reste indéfiniment en cours
d'émission, et les images doivent avoir été publiées une première fois (étape 1).
L'\textbf{application se fait dans l'ordre}, à périmètre restreint d'abord --- identité, secrets,
base, espace média --- puis complète, le plan étant lu intégralement à chaque fois
(\cref{subsec:incidents}, étape 4). La \textbf{chaîne de livraison dédiée} de l'étape 5 est recopiée
depuis celle de l'application pilote : cette duplication est l'écart le plus coûteux de la procédure,
et sa correction, un flux appelable commun, est identifiée. La quatrième contrainte est la plus
traître : les \textbf{cinq variables d'environnement} de l'étape 3 --- services collectés, chemins
d'authentification protégés, routage au bord, services supervisés, exclusions de règle --- doivent
être mises à jour \emph{ensemble}, chacune remplaçant un repli historique ; renseigner l'une sans y
réinclure les valeurs existantes coupe la collecte ou l'anti-force-brute de la plateforme en croyant
ajouter une application.

\begin{figure}[htbp]
\centering
\begin{tikzpicture}[
  x=1cm, y=1cm, font=\footnotesize,
  et/.style={nsysteme, text width=4.60cm, inner sep=3pt, minimum height=0.62cm},
  hc/.style={nabsent, text width=4.60cm, inner sep=3pt, minimum height=0.62cm},
  puce/.style={nabsent, text width=4.90cm, inner sep=3pt, minimum height=0.62cm,
               align=left}]

% ============ frontiere dans le code / hors du code =========================
\draw[frontiereplan] (7.25,0.95) -- (7.25,7.75);
\node[font=\footnotesize\bfseries, anchor=south] at (3.55,7.80) {dans le code};
\node[font=\footnotesize\bfseries, anchor=south] at (11.05,7.80) {hors du code};

% ============ les six etapes ================================================
\node[hc] (e1) at (3.55,7.25) {\textbf{1.} prérequis externes};
\node[et] (e2) at (3.55,6.50) {\textbf{2.} déclaration de l'application};
\node[hc] (e3) at (3.55,5.75) {\textbf{3.} cinq variables d'environnement};
\node[et] (e4a) at (3.55,5.00) {\textbf{4a.} application à périmètre restreint};
\node[nporte, text width=2.75cm] (pt) at (3.55,3.75) {porte : lecture intégrale du plan};
\node[et] (e4b) at (3.55,2.50) {\textbf{4b.} application complète};
\node[hc] (e5) at (3.55,1.75) {\textbf{5.} chaîne de livraison dédiée};
\node[hc] (e6) at (3.55,1.00) {\textbf{6.} correspondance de filtrage};
\foreach \a/\b in {e1/e2,e2/e3,e3/e4a,e4b/e5,e5/e6}{\draw[fluxdonnee] (\a) -- (\b);}
\draw[fluxdonnee] (e4a) -- (pt);
\draw[fluxcontrole] (pt) -- (e4b);

% ============ les cinq gestes hors du code ==================================
\node[puce] (p1) at (11.05,7.25) {enregistrement de domaine};
\node[puce] (p3) at (11.05,5.75) {fichier non versionné};
\node[puce] (p4) at (11.05,3.75) {dix ordres SQL, contrôlés quotidiennement};
\node[puce] (p5) at (11.05,1.75) {duplication de la chaîne};
\node[puce] (p6) at (11.05,1.00) {aucun test};
\draw[lienstructure] (e1.east) -- (p1.west);
\draw[lienstructure] (e3.east) -- (p3.west);
\draw[lienstructure] (pt.east) -- (p4.west);
\draw[lienstructure] (e5.east) -- (p5.west);
\draw[lienstructure] (e6.east) -- (p6.west);

% ============ legende interne ===============================================
\node[legendecadre, text width=15.00cm, anchor=north west, font=\scriptsize]
     at (0.05,0.45)
  {\legendetrait{fluxdonnee}{étape suivante}\hfill
   \legendetrait{fluxcontrole}{le plan est lu avant d'être appliqué}\hfill
   \legendetrait{lienstructure}{rattachement, jamais un flux}\\[2pt]
   losange : porte humaine\hfill
   cadre en tirets : \textbf{geste non descriptible en code}\hfill
   trait vertical épais : frontière du code};

\end{tikzpicture}
\caption[Accueil d'un locataire : six étapes et cinq gestes hors du code]{Accueil d'un locataire :
six étapes, et les cinq gestes qui restent hors du code. L'accueil est reproductible sur le plan de
l'identité et du déploiement ; les cinq gestes non descriptibles en code sont la mesure exacte de la
maturité multi-locataire du socle.}
\label{fig:accueil-locataire}
\end{figure}

\begin{alertbox}[boxsep=1pt, left=4pt, right=4pt, top=3pt, bottom=3pt]
S'ajoute à ces six étapes un geste non descriptible en code en l'état : \textbf{dix ordres SQL,
dont six révocations} qui retirent aux comptes applicatifs l'appartenance au rôle privilégié de
l'instance (annexe~\ref{ann:extraits-realisation}). Un contrôle quotidien vérifie qu'ils tiennent
toujours : \textbf{il détecte la dérive, il ne la prévient pas}. Le chemin par lequel cette
frontière se défait n'est pas la restauration elle-même mais la réconciliation qui la suit, et il
n'a aucune borne temporelle : le raisonnement complet, sa conséquence sur la conception et la
condition de clôture sont établis en \cref{subsec:isolation-donnees}.
\end{alertbox}

Le second locataire a été accueilli en recette le 07/08/2026 par instanciation du module
d'exécution paramétrique : compte de service dédié à deux rôles (client base de données, écriture
de journaux), base de données et utilisateur séparés, DNS et certificat gérés, six alertes et une
sonde de disponibilité ; depuis le 08/08/2026, les secrets et l'espace média sont chiffrés par clé
gérée dédiée. Le contrôle quotidien d'isolation du 19/08/2026 a produit \textbf{six vérifications
réussies sur six}, la connexion croisée entre bases étant effectivement refusée.

L'isolation obtenue porte sur quatre plans, dont un seul reste partagé, et cette délimitation vaut
mieux qu'une appréciation globale. \textbf{L'identité et les secrets} sont isolés sans réserve : le
compte de service de l'application hébergée ne détient \emph{aucun droit sur l'entrepôt de
supervision}, ce qui lui interdit structurellement de lire les données de sécurité, les siennes
comprises. \textbf{Les données} sont séparées logiquement --- base et utilisateur dédiés --- sur une
instance partagée, par les ordres décrits ci-dessus. \textbf{Le réseau est isolé depuis la migration
du 11 au \dateref} (\cref{subsec:migration-reseau}) : un sous-réseau par locataire, et quatre refus
nommés --- deux en sortie, deux en entrée --- interdisent chaque sens de traversée et journalisent
toute tentative sous un nom de règle exploitable comme signal de mouvement latéral. Le connecteur
unique, point de défaillance commun aux deux locataires, n'existe plus.

\textbf{Le plan qui reste partagé est celui du stockage.} Les deux locataires demeurent sur
\emph{une seule} instance de base de données et \emph{un seul} jeu de tables de supervision : la
ségrégation réseau empêche une identité d'atteindre les charges de l'autre locataire, non les deux
d'atteindre le même serveur, puisque les deux règles qui l'autorisent visent la même adresse ; le
contrôle reste au niveau du moteur de base. L'isolation des données est conçue, chiffrée et
ordonnancée en \cref{subsec:isolation-donnees}, et non appliquée --- la raison en est écrite plutôt
que tue : \textbf{deux migrations d'infrastructure dans la même quinzaine, conduites par un
opérateur unique, ne seraient pas soutenables}.

\begin{alertbox}[boxsep=1pt, left=4pt, right=4pt, top=3pt, bottom=3pt]
La portée exacte de cette frontière est celle établie en \cref{subsec:reseau-cible} : le sélecteur
disponible étant l'étiquette réseau et non le compte de service, elle est \textbf{opposable à une
charge de travail compromise, non à un opérateur capable de déployer}, qui poserait l'étiquette de
l'autre locataire. Ce second cas relève du resserrement du droit de déploiement
(\cref{subsec:delegation}), non du réseau ; et la conception ne tient pas au-delà de trois ou
quatre locataires, les refus croisés croissant comme le carré de leur nombre.
\end{alertbox}

Le seuil de bascule reste l'\textbf{accueil d'un troisième locataire} (\cref{sec:socle-service}),
mais son contenu a changé : sur les quatre travaux qu'il commandait, \textbf{deux sont faits} --- la
sortie par locataire, obtenue par un sous-réseau propre plutôt que par un connecteur dédié, et les
règles internes ciblées par identité, portées par la convention d'étiquetage ci-dessus. Les
\textbf{deux autres subsistent} : une clé de locataire obligatoire à l'ingestion et un cloisonnement
au stockage de l'entrepôt, auxquels s'ajoute l'instance de base dédiée au locataire porteur de donnée
irréparable (\cref{subsec:isolation-donnees}). Tous relèvent du plan de la donnée, aucun du réseau ;
le sous-réseau du troisième locataire est déjà créé et adressé, de sorte que son accueil est devenu,
côté réseau, une ligne de configuration.

L'expérimentation montre la réutilisabilité partielle du module d'exécution sans qualifier le socle
de produit multi-locataire abouti, et son coût de sécurité se dit symétriquement : le périmètre
d'usurpation accordé à l'identité de déploiement croît avec chaque application accueillie. Un dernier
piège appartient à la pire catégorie, celle des défaillances silencieuses : omettre le nouveau
service dans la liste des sources collectées rend l'application \emph{absente de l'entrepôt sans
produire la moindre erreur} --- l'application tourne, le socle la sert, et le centre d'opérations ne
la voit pas.

\subsection{La frontière de responsabilité}

Héberger l'application d'un tiers, fût-il du même groupe, oblige à trancher qui répond de quoi. Le
\cref{tab:frontiere-resp} énonce cette frontière sur quatorze domaines ; les cases où le socle
n'apporte rien y figurent au même titre que celles où il apporte tout, et un domaine reste laissé à
découvert par les deux parties --- c'est le plus utile à lire.

{\scriptsize
\setlength{\tabcolsep}{3pt}%
\begin{longtable}{|P{2.90cm}|P{6.25cm}|P{5.65cm}|}
\caption[Frontière de responsabilité entre le socle et l'éditeur]{Frontière de responsabilité
entre le socle et l'éditeur de l'application hébergée, état au \dateref.}
\label{tab:frontiere-resp}\\
\hline
\textbf{Domaine} & \textbf{Le socle (hébergeur)} & \textbf{L'éditeur de l'application hébergée} \\
\hline
\endfirsthead
\caption[]{Frontière de responsabilité entre le socle et l'éditeur (suite)} \\
\hline
\textbf{Domaine} & \textbf{Le socle (hébergeur)} & \textbf{L'éditeur de l'application hébergée} \\
\hline
\endhead
Noms de domaine, transit, certificats & Intégral : routage par hôte, certificat géré, politique
de chiffrement, transport strict & --- \\ \hline
Filtrage applicatif au bord & Intégral : politique unique, limitation de débit, blocage
géographique, motifs connus & Déclare ses chemins d'authentification à protéger \\ \hline
Réseau --- périmètre et sortie & Intégral : entrée par le seul répartiteur ; en sortie, refus par
défaut, cinq autorisations nommées par identité et par port, refus journalisés vers l'entrepôt.
\emph{Réserve} : les interfaces du fournisseur restent une destination autorisée et non inspectée
& --- \\ \hline
Réseau --- segmentation entre locataires & Intégral \emph{au plan réseau} depuis le 25/08/2026 : un
sous-réseau par locataire, refus croisés nommés et journalisés dans les deux sens, plus de
connecteur partagé (T20). \emph{Réserves} : opposable à une charge compromise, non à un opérateur
capable de déployer ; deux charges à sortie managée hors du pare-feu (T14) ; instance de base et
entrepôt partagés & Déclare ses charges, qui reçoivent l'étiquette de leur compte de service et le
sous-réseau de leur locataire \\ \hline
Exécution des conteneurs & Intégral : plateforme, quotas, sondes, mise à l'échelle, révisions,
entrée restreinte au répartiteur & Image conforme : exécution non privilégiée, port imposé, point
de santé sans dépendance externe \\ \hline
Identité machine et secrets & Intégral : compte de service dédié à moindre privilège, usurpation
ciblée, un secret par usage, clé gérée & Consomme les secrets par référence, ne les journalise
jamais \\ \hline
Base de données --- instance, sauvegarde, restauration & Intégral : instance régionale, adresse
privée, restauration à un instant donné ; clé gérée structurellement inapplicable, le champ étant
immuable après création. \emph{Réserve} : une restauration par clonage porte les deux bases &
Cohérence applicative après restauration, données hors base \\ \hline
Base de données --- schéma et isolation logique & Partiel : vérifie quotidiennement que la
séparation tient & Applique le durcissement ; geste non descriptible en code à ce jour \\ \hline
Stockage des médias et chiffrement au repos & Intégral au niveau de l'infrastructure : espace
dédié, accès public interdit, clé gérée & Sert les objets, gère les accès temporaires, et
\textbf{répond du chiffrement applicatif des données sensibles}, aujourd'hui absent \\ \hline
Code applicatif, logique métier, validation des entrées & \textbf{Aucune} sur la logique ; le
filtrage au bord n'arrête que des charges \emph{connues} & Intégral \\ \hline
Identités des utilisateurs finaux & \textbf{Aucune} : le second facteur du socle ne couvre que les
comptes du tableau de bord & Intégral : authentification, autorisation et second facteur des
utilisateurs de l'application \\ \hline
Journalisation, détection et notification & Journalisation d'infrastructure, normalisation,
conservation, R1--R7, rattachement ATT\&CK, enrichissement. \emph{Réserve} : la journalisation
métier est ingérée sans être exploitée, et l'éditeur n'est pas notifié & Émet ses journaux, déclare
ses exclusions de règle, conserve son alerte fonctionnelle \\ \hline
Chaîne de livraison et analyse d'image & Intégral : authentification sans clé, registre unique,
analyse de vulnérabilité, restitution à l'entrepôt & Corrige les vulnérabilités remontées \\ \hline
Conformité, base légale, droits des personnes & \textbf{Aucune} : des moyens techniques, jamais
une gouvernance & Intégral : base légale, information des personnes, exercice des droits, durées de
conservation \\
\hline
\end{longtable}
}

\begin{keybox}[boxsep=1pt, left=4pt, right=4pt, top=3pt, bottom=3pt]
Le socle sécurise le bord, le réseau, l'exécution, l'identité machine et la chaîne de livraison.
La logique métier, les identités des utilisateurs finaux et la conformité restent à la charge de
l'éditeur. Cette frontière n'est pas un partage de commodité : elle correspond à ce qu'un
hébergeur peut techniquement contrôler.
\end{keybox}

Une ligne appelle une lecture particulière, car elle est vide des deux côtés : le chiffrement
applicatif des données sensibles. Le socle chiffre les supports, il ne chiffre pas les champs, et
l'éditeur ne l'a pas fait. Ce vide n'est pas un oubli de rédaction ; c'est l'exigence qu'un contrat
d'hébergement devrait porter en premier, et le mémoire la nomme plutôt que de la répartir
arbitrairement. La segmentation réseau entre locataires occupait cette même case jusqu'à la
migration du 11 au \dateref{} (\cref{subsec:migration-reseau}) : elle a basculé côté hébergeur, seul
domaine du tableau à avoir changé de camp. \textbf{Un domaine sort du découvert lorsqu'une partie
l'assume dans son code, non lorsqu'un contrat le lui attribue.}

\section{Conduite du projet}

J'ai conduit seul toutes les étapes du projet : audit du dépôt existant, rédaction des exigences,
conception de l'architecture, écriture des modules d'infrastructure, construction des deux chaînes,
traduction des règles de détection, mise en place de la couche d'enrichissement et du tableau de
bord, puis exploitation quotidienne du socle en recette --- celle qui a produit les cinq incidents
ci-dessous. Un assistant de programmation a servi à produire des squelettes répétitifs, des tests et
des relectures, jamais à décider d'une option d'architecture ou de sécurité, et aucune production
assistée n'a été retenue sans exécution : elle est passée par les portes de la chaîne de livraison
puis par une vérification en conditions réelles. Les deux changements de périmètre significatifs ---
abandon du quatrième projet cloud d'amorçage (\cref{sec:env-travail}) et fédération d'identité
simplifiée (\cref{subsec:wif}) --- ont été documentés au moment de leur décision.

\subsection{Difficultés rencontrées}
\label{subsec:incidents}

Chacun des cinq incidents ci-dessous a produit une correction et une leçon transférable.

\textbf{Indisponibilité liée au démarrage à froid.} Le service d'encodage, réduit à zéro instance
hors traitement, doit démarrer et charger le modèle lorsqu'une requête arrive après une période
d'inactivité. Les mesures rapportées début août ont établi une durée moyenne d'environ vingt-sept
secondes, avec un pic à quatre-vingt-quatorze secondes (\cref{graph:coldstart}) --- un tout autre
ordre de grandeur que l'estimation initiale, conservée avec la mesure qui la contredit plutôt que
corrigée discrètement. Le budget de la sonde de disponibilité a été porté au-dessus du pic observé ;
maintenir une instance active en permanence aurait supprimé le problème, pour un coût
disproportionné.

\textbf{Blocage collectif par la limitation de débit.} Le mécanisme initial, un seuil unique par
adresse IP, bloquait des utilisateurs légitimes : le tableau de bord et l'API partageant une
adresse de sortie commune, l'activité normale de quelques analystes suffisait à franchir le seuil.
La correction a introduit la limitation à deux niveaux décrite au \cref{ch:conception}, complétée côté
applicatif par le compteur indexé sur le défi (\cref{subsec:authent-dashboard}).

\textbf{État divergent après une application partiellement échouée.} Le 08/08/2026, une modification
d'infrastructure appliquée sur la recette s'est terminée en échec partiel : une contrainte du
fournisseur, non évaluée à la planification, a refusé une partie des changements. L'état local a
pourtant enregistré une configuration jamais appliquée, si bien que la planification suivante ne
signalait plus aucune différence alors que l'infrastructure réelle en présentait une. L'écart a été
détecté par comparaison explicite entre l'état enregistré et les valeurs lues chez le fournisseur,
puis corrigé par resynchronisation. Le constat à retenir est factuel : un échec partiel peut laisser
l'état local différent de l'état réel, et la planification suivante hérite de cette erreur.

\textbf{Porte SAST en faux vert.} L'action Semgrep utilisée par la chaîne échouait en interne sur
une incompatibilité de sévérité, tandis que l'étape apparaissait verte : le 16/08/2026, toutes les
tâches de la chaîne étaient en succès alors qu'aucune règle n'avait été évaluée malgré la porte
annoncée. Le 19/08/2026, l'action a été remplacée par un appel direct à l'outil en ligne de
commande, avec propagation du code de sortie ; la première exécution réelle, numéro 32264602112, a
analysé 439 fichiers avec 612 règles, produit 74 constats et terminé avec le code 1. La chaîne en
échec prouve ici le rétablissement du contrôle ; elle ne prouve pas la remédiation des 74 constats,
dont le triage reste nécessaire.

\begin{figure}[!htbp]
\centering
\IfFileExists{figures/K23.png}{%
  \setlength{\fboxrule}{0.4pt}\setlength{\fboxsep}{0pt}%
  \fcolorbox{black!30}{white}{\includegraphics[width=0.86\textwidth]{figures/K23.png}}%
}{%
  \begin{minipage}{0.86\textwidth}
  \begin{extraitconf}[breakable=false, fontupper=\footnotesize,
      title={Capture K23 : spécification de prise}]
  \textbf{Vue :} historique des exécutions de la chaîne de livraison applicative, filtré sur le flux
  qui porte la porte d'analyse statique, cadré sur deux lignes voisines et l'en-tête de colonnes.
  \textbf{Dates portées par la capture :} 16/08/2026 et 19/08/2026, preuve d'archive.\par\smallskip
  \textbf{Contenu probant :} la seule chose que ce chapitre ne peut pas écrire ---
  \emph{qu'une porte inerte a l'air de fonctionner}. Ligne du 16/08 : toutes les tâches en succès,
  alors qu'aucune règle n'a été évaluée. Ligne du 19/08 : la même chaîne en échec après remise en
  service, exécution 32264602112.\par\smallskip
  \textbf{Masquage :} nom du dépôt et de l'organisation, courriels des auteurs de
  \emph{commit}.
  \end{extraitconf}
  \end{minipage}%
}
\caption[Porte d'analyse statique : le faux vert du 16/08 et l'échec du 19/08]{Historique des
exécutions de la chaîne de livraison applicative. À y lire, sur deux lignes voisines : le
16/08/2026, toutes les tâches en succès alors que la porte d'analyse statique n'évaluait aucune
règle ; le 19/08/2026, la même chaîne en échec après remise en service. Le vert de la première
ligne est l'objet de la preuve : c'est lui, et non un journal, qui montre qu'un contrôle inerte est
indiscernable d'un contrôle qui passe.}
\label{fig:capture-K23}
\end{figure}

\textbf{Une règle de sortie qui n'aurait rien gouverné.} Le refus par défaut en sortie figurait au
plan de correction bien avant d'être posé, et la lecture croisée de la topologie et du paramétrage
des charges a montré qu'il aurait été \emph{sans effet}, pour la raison analysée en
\cref{subsec:reseau-cible} : le trafic ne passait pas là où la règle se serait appliquée. La
correction a donc consisté à déplacer d'abord le trafic, puis à l'observer sept jours, et à
n'interdire qu'ensuite (\cref{subsec:migration-reseau}). La leçon est transférable telle quelle :
\textbf{un contrôle qui ne voit pas passer le trafic n'est pas un contrôle, c'est une
déclaration} --- avant de mesurer l'efficacité d'une règle, il faut établir qu'elle est sur le
chemin.

\subsection{Planning de travail}
\label{sec:planning-travail}

La période de stage s'étend du \datedebutstage{} au \datefinstage{}, et la \cref{fig:gantt} la couvre
intégralement. Faute de journal quotidien avant le 29/07/2026, ses barres n'affichent pas une fausse
précision au jour près ; les losanges M1 à M7 renvoient à des jalons datés par une trace
technique.

% Placement [p] : la norme ESPRIT demande que le planning occupe une page
% complète, en fin de description du travail. Le spécificateur de page de
% flottant est ce qui le garantit, sans laisser de page blanche avant lui.
\begin{figure}[p]
\centering
\begin{tikzpicture}[x=0.0625cm, y=0.72cm, font=\footnotesize,
  phase/.style={anchor=east, xshift=-4mm, text width=3.15cm, align=right,
    font=\scriptsize},
  % Trois teintes seulement, correspondant aux trois familles de phases :
  % cadrage, construction, validation et cloture. Le degrade code une famille.
  analyse/.style={draw=menalencre, thick, rounded corners=1pt, fill=menalplanA},
  conception/.style={draw=menalencre, thick, rounded corners=1pt, fill=menalplanA},
  realisation/.style={draw=menalencre, thick, rounded corners=1pt, fill=menalplanB},
  validation/.style={draw=menalencre, thick, rounded corners=1pt, fill=menalplanC},
  documentation/.style={draw=menalencre, thick, rounded corners=1pt, fill=menalplanC},
  mois/.style={font=\scriptsize\bfseries},
  jalon/.style={diamond, aspect=1.1, draw=menalencre, fill=menalplanC, thick,
    inner sep=0pt, minimum width=2.4mm, minimum height=2.4mm},
  repere/.style={anchor=north, font=\scriptsize}
]
  % En-tête mensuel : abscisse exprimée en jours depuis le 03/03/2026.
  % Les deux derniers mois sont fusionnés en une cellule pour éviter que
  % « sept. » ne déborde d'une cellule de six jours.
  \foreach \a/\b/\nom in {0/29/mars,29/59/avril,59/90/mai,90/120/juin,
    120/151/juillet,151/188/{août -- sept.}}{
    \draw[fill=black!5, draw=black!35] (\a,8.75) rectangle (\b,9.35);
    \node[mois] at ({(\a+\b)/2},9.05) {\nom};
    \draw[black!18] (\a,0.55) -- (\a,8.75);
  }
  \draw[black!18] (188,0.55) -- (188,8.75);

  % Les sept verticales de jalon ont ete retirees : sept traits coupant huit
  % barres formaient une trame parasite sans rien ajouter, les jalons etant
  % portes par les losanges de l'axe et dates un a un dans la legende.

  \node[phase] at (0,8.15) {PH0 --- Immersion\\et cadrage};
  \draw[analyse] (0,7.83) rectangle (24,8.47);
  \node[font=\scriptsize, anchor=west] at (26,8.15) {03/03 -- 27/03};

  \node[phase] at (0,7.15) {PH1 --- Étude\\de l'existant};
  \draw[analyse] (13,6.83) rectangle (148,7.47);
  \node[font=\scriptsize] at (80.5,7.15) {analyse progressive jusqu'à l'audit figé};

  \node[phase] at (0,6.15) {PH2 --- Risques\\et exigences};
  \draw[conception] (59,5.83) rectangle (150,6.47);
  \node[font=\scriptsize] at (104.5,6.15) {mai -- juillet};

  \node[phase] at (0,5.15) {PH3 --- Architecture\\Zero Trust};
  \draw[conception] (90,4.83) rectangle (153,5.47);
  \node[font=\scriptsize] at (121.5,5.15) {juin -- début août};

  \node[phase] at (0,4.15) {PH4 --- IaC et\\chaînes de livraison};
  \draw[realisation] (120,3.83) rectangle (168,4.47);
  \node[font=\scriptsize] at (144,4.15) {juillet -- 18/08};

  \node[phase] at (0,3.15) {PH5 --- Détection\\et enrichissement};
  \draw[realisation] (139,2.83) rectangle (175,3.47);
  \node[font=\scriptsize] at (157,3.15) {20/07 -- 25/08};

  \node[phase] at (0,2.15) {PH6 --- Validation\\et corrections};
  \draw[validation] (153,1.83) rectangle (182,2.47);
  \node[font=\scriptsize] at (167.5,2.15) {03/08 -- 31/08};

  \node[phase] at (0,1.15) {PH7 --- Rapport,\\transfert et clôture};
  \draw[documentation] (168,0.83) rectangle (188,1.47);
  \node[font=\scriptsize, anchor=east] at (166,1.15) {18/08 -- 07/09};

  % Jalons exactement datés : losanges posés sur l'axe, sans étiquette en
  % place — les jalons couvrent neuf jours et sont regroupés sous l'axe,
  % chaque date étant donnée dans la légende ci-dessous.
  \foreach \x in {148,151,153,157,169,175,188}{\node[jalon] at (\x,0.55) {};}
  \node[repere] at (152,0.28) {M1--M4};
  \node[repere] at (172,0.28) {M5--M6};
  \node[repere, anchor=north east] at (190,0.28) {M7};
\end{tikzpicture}

\vspace{4pt}
\begin{minipage}[t]{0.48\textwidth}\scriptsize
\textbf{M1 --- 29/07 :} audit statique de l'existant figé.\\
\textbf{M2 --- 01/08 :} export int8 rejeté, fp32 conservé.\\
\textbf{M3 --- 03/08 :} restauration chronométrée.\\
\textbf{M4 --- 07/08 :} accueil du locataire ELSON.
\end{minipage}\hfill
\begin{minipage}[t]{0.48\textwidth}\scriptsize
\textbf{M5 --- 19/08 :} correctifs et scénario de bout en bout.\\
\textbf{M6 --- 25/08 :} date de référence et consolidation des preuves.\\
\textbf{M7 --- 07/09 :} fin du stage et transfert des livrables.
\end{minipage}
\caption[Planning consolidé du stage, du \datedebutstage{} au \datefinstage{}]{Planning consolidé
du stage, du \datedebutstage{} au \datefinstage{}. Les trois teintes de barre distinguent le
cadrage (PH0--PH3), la construction (PH4--PH5) et la validation et la clôture (PH6--PH7). Les
barres sont une reconstruction macroscopique ; les losanges signalent les jalons exactement datés,
dont la légende donne la date une à une.}
\label{fig:gantt}
\end{figure}


\section*{Conclusion du chapitre}
\addcontentsline{toc}{section}{Conclusion du chapitre}

L'infrastructure est décrite en code, et la frontière de ce principe est énoncée par six exceptions
nommées plutôt que par une revendication d'exhaustivité. La chaîne de livraison compte trois portes
bloquantes, cinq familles de contrôles réalisées et trois arbitrages assortis de leur condition de
réévaluation ; une seule famille voit son résultat quitter la chaîne, ce qui rend la boucle F6
possible. La chaîne de provisionnement s'arrête volontairement avant l'application, dont
l'automatisation exigerait une concentration de privilège que le \cref{ch:conception} refuse, et le tableau de
bord ne détient aucune autorisation propre. L'accueil du second locataire isole l'identité, les
secrets, les données et, depuis la migration du 11 au \dateref, le réseau ; le stockage reste
partagé, faute qu'une seconde migration soit soutenable dans la même quinzaine par un opérateur
unique. Deux limites restent assumées : deux charges conservent une sortie managée hors du pare-feu
(T14, non conforme), et \textbf{l'enrichissement n'alimente pas le score d'incident}.
